\scanpage{13}
% archive foliation: 12
la connaissance des $\mu_i$.

Si on se donne les \inserted{entiers} $\nu_j\geq0$
\deleted{\unsure{$=\operatorname{rg}\pi_j(L_*)$ $(j\geq0)$}}
$(j\geq0)$ et une fonction \inserted{polynomiale} $f(n)$ à valeurs
entières définie pour $n$ grand, à quelles
conditions existe-t-il $L_*$ str. spécial
avec $\nu_j=\operatorname{rg}\pi_j(L_*)$ $(j\geq0)$ et
$f(n)=f_{L_*}(n)$ pour $n$ grand? Notons
que $g=f-\sum_{j\geq0}\nu_jP_j$ est encore une
fonction \deleted{entière} \inserted{polynomiale} à valeurs entières,
donc de la forme $\sum_{i\geq0}\mu_i\binom{n+1}{i}$
où $(\mu_i)\in\mathbf Z^{(\mathbf N)}$ est uniquement déterminé.
Alors la condition cherchée est
\[
\mu_0=0,\quad\mu_i\geq0\quad\text{si }i\geq1.
\]
Notons que \deleted{\unsure{$\mu_0$}}
\[
\mu_0=g(-1)=f(-1)-\sum_{j\geq0}\nu_j
\frac{(-1)(-2)\cdots(-j)}{j!}
\]
donc on a aussi
\[
(\mu_0=0)\Longleftrightarrow
\left(f(-1)=\underset{\chi((\nu_j))}{\underbrace{\sum(-1)^j\nu_j}}\right)
\]

\deleted{\unsure{Déf.}
\deleted{Complexes}
$k$-Modules $\frac12$ simpliciaux $L_*$ \deleted{\unsure{parfaits}}
``\inserted{str.}\unsure{chaîne-parfait}''
\inserted{(strict. \unsure{chaîne-parfait})}
$=$ tels que $L^!_*$ soit
un complexe de chaînes de $k$-modules
\deleted{\unsure{str.}}
parfait (str. parfait) $=$}
